% Two forgetting models: what is retained after three probes of length tau = 4,
% each followed by an interruption (gamma = 1/2, retained progress starts at 0).
%   this talk: R' = gamma (R + tau)  -> every interruption shrinks ALL progress:
%              probe 1 keeps gamma^3 tau = 0.5, probe 2 gamma^2 tau = 1, probe 3 gamma tau = 2
%   variant:   R' = R + gamma tau    -> each probe shrinks only once: 2, 2, 2
% x = amount of work (1 unit = 1 time unit of processing); colour = probe.
\begin{tikzpicture}[x=0.8cm,y=1cm]
  % --- row y=2: the work that was done -------------------------------------
  \node[rowlabel] at (0,2.3) {work done};
  \job[2]{jobA}{0}{4}{probe 1}
  \job[2]{jobB}{4}{8}{probe 2}
  \job[2]{jobC}{8}{12}{probe 3}
  \node[endlabel] at (12,2.3) {$\tau = 4$, \ $\gamma = \tfrac12$};
  % --- row y=1: this talk, all progress shrinks at every interruption ------
  \node[rowlabel] at (0,1.3) {this talk};
  \job[1]{jobA}{0}{0.5}{$\tfrac12$}
  \job[1]{jobB}{0.5}{1.5}{$1$}
  \job[1]{jobC}{1.5}{3.5}{$2$}
  \node[endlabel] at (3.5,1.3) {retained: $3.5$ \ \ (old work keeps fading)};
  % --- row y=0: variant, only the last probe shrinks, once -----------------
  \node[rowlabel] at (0,0.3) {variant};
  \job{jobA}{0}{2}{$2$}
  \job{jobB}{2}{4}{$2$}
  \job{jobC}{4}{6}{$2$}
  \node[endlabel] at (6,0.3) {retained: $6$ \ \ (each probe fades once)};
  % --- axis: amount of work --------------------------------------------------
  \draw[axis] (0,-0.15) -- (12.6,-0.15) node[right,lbl]{work};
  \foreach \t in {0,2,...,12}
    \draw[tick] (\t,-0.15) -- ++(0,-0.1) node[below,lbl]{$\t$};
\end{tikzpicture}
